\section{Implementation and Numerical Qualification}
\label{sec:implementation}

\subsection{Arithmetic paths after shared preparation}
The implementation extends a llama.cpp-derived decoder and reuses existing quantized kernels~\cite{llama_cpp_fewrow,loop4code}. \Cfour{} selects admitted few-column paths for two to four queries; \Pfour{} organizes independent four-column groups in wider prefill; \Local{} places physical-view refinement inside attention. These choices expose different execution freedoms.

Table~\ref{tab:dispatch} distinguishes the RTX paths after the same Q8\_1 activation preparation. B is the controlled scalar-loop baseline, C1 groups independent scalar channel/query work into one dispatch, and C4 admits the inherited few-column specialization. The initial and confirmation studies share their numerical core, worker binaries and models.

\papertable{t01_dispatch}{Target-projection dispatch in the independent RTX B/C1/C4 studies. All three paths begin after the same Q8\_1 activation preparation; counts refer to one examined MMVQ operation.}{tab:dispatch}

RTX C4 admits Q8\_0 weights on the GENERIC device table, registered $K\geq1024$ block alignment, even output dimension $M$, and $2\leq q\leq4$. Fusion, FlashAttention, CUDA graphs and programmatic dependent launch are disabled in that study. The A5000 C1/C4 bridge uses its own q2--4 selector for Q4\_K/Q6\_K paths; C1 retains the scalar-channel representation and both arms have P4 off. The artifact source map binds these distinct builds to their executed routines.

\subsection{Preserve arithmetic within groups; organize between them}
For $q=4g$ known columns, P4 interprets supported contiguous activation columns as $g$ independent four-column groups after shared preparation. Each retains its admitted arithmetic; no reduction crosses groups. The A5000 graph keeps input/output shapes $[K,1,q]$ and $[M,1,q]$, and the MMVQ dispatcher supplies the group interpretation at launch. The wide predicate is $8\leq q\leq32$, \mbox{$q\bmod4=0$}; other shapes retain their registered paths. The request study exercises q32 packed prefill, while separate q16/q24 panels qualify additional numerical cases.

The LOCAL panel uses C4 with P4 disabled and reuses the frozen A5000 prefill numerical core. At its crossing, SIMPLE's q1 fragment partly uses gate/up fusion, while LOCAL's q4 path is unfused. This behavior differs from the fusion-disabled RTX bridge and remains part of the measured LOCAL treatment.

\subsection{Qualify the complete numerical interface}
Qualification checks complete requested vocabulary vectors, active KV in canonical logical order, and successor state. The completed comparisons in Table~\ref{tab:qualification} match their bound references: vectors and active KV agree bitwise, and consumer decisions and successor state agree. Heavy captures are separate from clean timing. Conditional proposal rows use a scalar reference with the same supplied prefix; autonomous continuation checks the emitted sequence and stopping behavior.

\papertable{t02_qualification}{Numerical qualification of the bound CUDA paths. These are original-machine complete-object checks; Table~\ref{tab:access} distinguishes the independently readable public subset.}{tab:qualification}

Sentinel captures test full objects, while timed requests retain endpoint and actual-work checks. LOCAL compares owned rows, tentative and committed state, and continuation from each arm's own commit. RTX confirmation repeats pristine-reference captures on two sentinel roots and checks timed-request state/frontier identities, pending and transaction state, work, dispatch, and graph/view identity. Sections~\ref{supp:access} and~\ref{supp:rtx-confirmation} retain object counts and access scope. Structural admission guards enforce the dense-cache, greedy domain in Section~\ref{sec:reference}; complete-vector and state comparisons test the exercised numerical instances.
